\thispagestyle{plain}
\begin{center}
{\color{gilt}\pgfornament[width=5cm]{87}}\par\vspace{4pt}
{\fontsize{24}{28}\selectfont\color{navy}\scshape Contents}\par\vspace{4pt}
{\color{gilt}\pgfornament[width=5cm,symmetry=h]{87}}
\end{center}
\vspace{0.5cm}
\newcommand{\tocrow}[3]{\noindent{\large\color{navy}\scshape NEET #1 Physics}\dotfill{\large\color{crimson}\pageref{#2}}\par\vspace{2pt}{\small\itshape\hspace*{1em}#3}\par\vspace{6pt}}
{\leftskip=2cm\rightskip=2cm
\tocrow{2026}{p2026}{Re-examination, 21 June 2026 \textperiodcentered{} Test Booklet Code 50}
\tocrow{2025}{p2025}{4 May 2025 \textperiodcentered{} Test Booklet Code 45}
\tocrow{2024}{p2024}{5 May 2024 \textperiodcentered{} Test Booklet Code S6}
\tocrow{2023}{p2023}{7 May 2023 \textperiodcentered{} Test Booklet Code G6}
\tocrow{2022}{p2022}{17 July 2022 \textperiodcentered{} Test Booklet Code R6}
\tocrow{2021}{p2021}{12 September 2021 \textperiodcentered{} Test Booklet Code M4}
\tocrow{2020}{p2020}{13 September 2020 \textperiodcentered{} Test Booklet Code E1 \textperiodcentered{} Q 136--180}
\tocrow{2019}{p2019}{5 May 2019 \textperiodcentered{} Test Booklet Code P1}
\tocrow{2018}{p2018}{6 May 2018 \textperiodcentered{} Test Booklet Code AA}
\tocrow{2017}{p2017}{7 May 2017 \textperiodcentered{} Test Booklet Code B \textperiodcentered{} Q 136--180}
\tocrow{2016}{p2016}{Phase I, 1 May 2016 \textperiodcentered{} Test Booklet Code A}
\noindent{\large\color{navy}\scshape Eleven-Year Summary}\dotfill{\large\color{crimson}\pageref{psummary}}\par\vspace{2pt}{\small\itshape\hspace*{1em}Direct, applied and deleted-portion counts, 2016--2026}\par\vspace{6pt}
\noindent{\large\color{navy}\scshape The 15 Most Repeated Topics}\dotfill{\large\color{crimson}\pageref{ptopics}}\par\vspace{2pt}{\small\itshape\hspace*{1em}Topic counts over eleven years, and questions asked again}\par\vspace{6pt}
\noindent{\large\color{navy}\scshape How to Prepare for NEET Physics}\dotfill{\large\color{crimson}\pageref{pplan}}\par\vspace{2pt}{\small\itshape\hspace*{1em}Sources, method and plan of action from this analysis}\par\vspace{6pt}
\par}
\vfill
\begin{center}
\begin{minipage}{0.78\textwidth}\centering\small
\textbf{\color{navy}About this book.} The book holds the Physics section of the main NEET (UG) paper of each year from 2016 to 2026: 45 questions (2016--2020, 2025, 2026) or 50 questions (2021--2024, Section A 35 and Section B 15, of which any 10 were to be attempted), four marks each with one mark deducted for a wrong answer. For 2026 the re-examination of 21 June is used, because the 3 May examination was cancelled. Questions keep their original numbers and option labels, so they match the official answer key of the test booklet code named above. Every question has been retyped and every diagram redrawn from the English version of the paper. Every question was checked one by one against the four 2026-27 NCERT books (Class XI Part 1 \& 2, Class XII Part 1 \& 2) and carries a tag: \textsc{Direct} (answer stated in the textbook, an NCERT solved Example or Exercise, or a single printed formula), \textsc{Applied} (NCERT concept used in several steps or a new situation), or \textsc{Deleted Portion} (topic removed from the rationalised syllabus; not analysed). Direct and Applied tags give the NCERT page with a tappable link to that chapter's official PDF. Each paper is followed by its answer key (the official NTA final key for 2021, 2022, 2023, 2025 and 2026, checked by the editor; solved by the editor for the other years) and a one-page analysis. The book ends with an eleven-year summary, the 15 most repeated topics and a preparation plan.\\[8pt]
{\color{gilt}\pgfornament[width=2.5cm]{15}}\\[4pt]
{\itshape Prepared by Sumeet Birajadar for Science Monk online classes.}
\end{minipage}
\end{center}
\clearpage
